\documentclass[11pt]{article}
\newcommand{\ListaTitulo}{Gabarito --- Lista 05}
% Preâmbulo compartilhado — MATD48, Listas de Exercícios 2026
% Usado por Lista{NN}.tex e Gabarito{NN}.tex via % Preâmbulo compartilhado — MATD48, Listas de Exercícios 2026
% Usado por Lista{NN}.tex e Gabarito{NN}.tex via % Preâmbulo compartilhado — MATD48, Listas de Exercícios 2026
% Usado por Lista{NN}.tex e Gabarito{NN}.tex via \input{preamble}

\usepackage[utf8]{inputenc}
\usepackage[T1]{fontenc}
\usepackage[brazil]{babel}
\usepackage{fullpage}
\usepackage{amsmath,amssymb}
\usepackage{enumitem}
\usepackage{xcolor}
\usepackage{fancyhdr}
\usepackage{titlesec}
\usepackage{listings}
\usepackage{hyperref}
\usepackage{booktabs}
\usepackage{graphicx}

\definecolor{matd48azul}{HTML}{0B4F6C}
\definecolor{matd48verde}{HTML}{2E8B57}

\titleformat{\section}{\normalfont\Large\bfseries\color{matd48azul}}{\thesection}{1em}{}
\titleformat{\subsection}{\normalfont\large\bfseries\color{matd48azul}}{\thesubsection}{1em}{}

\providecommand{\ListaTitulo}{Lista de Exercícios}

\setlength{\headheight}{14pt}
\pagestyle{fancy}
\fancyhf{}
\lhead{\small MATD48 --- Planejamento de Experimentos A}
\rhead{\small \ListaTitulo}
\cfoot{\thepage}
\renewcommand{\headrulewidth}{0.4pt}

\lstset{
  basicstyle=\ttfamily\small,
  breaklines=true,
  frame=single,
  columns=fullflexible,
  backgroundcolor=\color{gray!8},
  commentstyle=\color{matd48verde},
  keywordstyle=\color{matd48azul}\bfseries,
  language=R,
  extendedchars=true,
  literate=%
    {á}{{\'a}}1 {à}{{\`a}}1 {â}{{\^a}}1 {ã}{{\~a}}1
    {é}{{\'e}}1 {ê}{{\^e}}1 {í}{{\'i}}1 {ó}{{\'o}}1
    {ô}{{\^o}}1 {õ}{{\~o}}1 {ú}{{\'u}}1 {ç}{{\c c}}1
    {Á}{{\'A}}1 {À}{{\`A}}1 {Â}{{\^A}}1 {Ã}{{\~A}}1
    {É}{{\'E}}1 {Ê}{{\^E}}1 {Í}{{\'I}}1 {Ó}{{\'O}}1
    {Ô}{{\^O}}1 {Õ}{{\~O}}1 {Ú}{{\'U}}1 {Ç}{{\c C}}1
}

\hypersetup{colorlinks=true, linkcolor=matd48azul, urlcolor=matd48azul, citecolor=matd48azul}

\newcommand{\cabecalho}[2]{%
  \begin{center}
    {\Large\bfseries\color{matd48azul} #1}\\[0.3em]
    {\large #2}\\[0.3em]
    \rule{\linewidth}{0.6pt}
  \end{center}
  \vspace{0.5em}
}

\newenvironment{questao}[1]{\vspace{1em}\noindent\textbf{Questão #1.}\hspace{0.4em}}{\par}
\newenvironment{solucao}{\vspace{0.3em}\noindent\textit{Solução.}\hspace{0.4em}\color{black!85}}{\par\vspace{0.5em}}


\usepackage[utf8]{inputenc}
\usepackage[T1]{fontenc}
\usepackage[brazil]{babel}
\usepackage{fullpage}
\usepackage{amsmath,amssymb}
\usepackage{enumitem}
\usepackage{xcolor}
\usepackage{fancyhdr}
\usepackage{titlesec}
\usepackage{listings}
\usepackage{hyperref}
\usepackage{booktabs}
\usepackage{graphicx}

\definecolor{matd48azul}{HTML}{0B4F6C}
\definecolor{matd48verde}{HTML}{2E8B57}

\titleformat{\section}{\normalfont\Large\bfseries\color{matd48azul}}{\thesection}{1em}{}
\titleformat{\subsection}{\normalfont\large\bfseries\color{matd48azul}}{\thesubsection}{1em}{}

\providecommand{\ListaTitulo}{Lista de Exercícios}

\setlength{\headheight}{14pt}
\pagestyle{fancy}
\fancyhf{}
\lhead{\small MATD48 --- Planejamento de Experimentos A}
\rhead{\small \ListaTitulo}
\cfoot{\thepage}
\renewcommand{\headrulewidth}{0.4pt}

\lstset{
  basicstyle=\ttfamily\small,
  breaklines=true,
  frame=single,
  columns=fullflexible,
  backgroundcolor=\color{gray!8},
  commentstyle=\color{matd48verde},
  keywordstyle=\color{matd48azul}\bfseries,
  language=R,
  extendedchars=true,
  literate=%
    {á}{{\'a}}1 {à}{{\`a}}1 {â}{{\^a}}1 {ã}{{\~a}}1
    {é}{{\'e}}1 {ê}{{\^e}}1 {í}{{\'i}}1 {ó}{{\'o}}1
    {ô}{{\^o}}1 {õ}{{\~o}}1 {ú}{{\'u}}1 {ç}{{\c c}}1
    {Á}{{\'A}}1 {À}{{\`A}}1 {Â}{{\^A}}1 {Ã}{{\~A}}1
    {É}{{\'E}}1 {Ê}{{\^E}}1 {Í}{{\'I}}1 {Ó}{{\'O}}1
    {Ô}{{\^O}}1 {Õ}{{\~O}}1 {Ú}{{\'U}}1 {Ç}{{\c C}}1
}

\hypersetup{colorlinks=true, linkcolor=matd48azul, urlcolor=matd48azul, citecolor=matd48azul}

\newcommand{\cabecalho}[2]{%
  \begin{center}
    {\Large\bfseries\color{matd48azul} #1}\\[0.3em]
    {\large #2}\\[0.3em]
    \rule{\linewidth}{0.6pt}
  \end{center}
  \vspace{0.5em}
}

\newenvironment{questao}[1]{\vspace{1em}\noindent\textbf{Questão #1.}\hspace{0.4em}}{\par}
\newenvironment{solucao}{\vspace{0.3em}\noindent\textit{Solução.}\hspace{0.4em}\color{black!85}}{\par\vspace{0.5em}}


\usepackage[utf8]{inputenc}
\usepackage[T1]{fontenc}
\usepackage[brazil]{babel}
\usepackage{fullpage}
\usepackage{amsmath,amssymb}
\usepackage{enumitem}
\usepackage{xcolor}
\usepackage{fancyhdr}
\usepackage{titlesec}
\usepackage{listings}
\usepackage{hyperref}
\usepackage{booktabs}
\usepackage{graphicx}

\definecolor{matd48azul}{HTML}{0B4F6C}
\definecolor{matd48verde}{HTML}{2E8B57}

\titleformat{\section}{\normalfont\Large\bfseries\color{matd48azul}}{\thesection}{1em}{}
\titleformat{\subsection}{\normalfont\large\bfseries\color{matd48azul}}{\thesubsection}{1em}{}

\providecommand{\ListaTitulo}{Lista de Exercícios}

\setlength{\headheight}{14pt}
\pagestyle{fancy}
\fancyhf{}
\lhead{\small MATD48 --- Planejamento de Experimentos A}
\rhead{\small \ListaTitulo}
\cfoot{\thepage}
\renewcommand{\headrulewidth}{0.4pt}

\lstset{
  basicstyle=\ttfamily\small,
  breaklines=true,
  frame=single,
  columns=fullflexible,
  backgroundcolor=\color{gray!8},
  commentstyle=\color{matd48verde},
  keywordstyle=\color{matd48azul}\bfseries,
  language=R,
  extendedchars=true,
  literate=%
    {á}{{\'a}}1 {à}{{\`a}}1 {â}{{\^a}}1 {ã}{{\~a}}1
    {é}{{\'e}}1 {ê}{{\^e}}1 {í}{{\'i}}1 {ó}{{\'o}}1
    {ô}{{\^o}}1 {õ}{{\~o}}1 {ú}{{\'u}}1 {ç}{{\c c}}1
    {Á}{{\'A}}1 {À}{{\`A}}1 {Â}{{\^A}}1 {Ã}{{\~A}}1
    {É}{{\'E}}1 {Ê}{{\^E}}1 {Í}{{\'I}}1 {Ó}{{\'O}}1
    {Ô}{{\^O}}1 {Õ}{{\~O}}1 {Ú}{{\'U}}1 {Ç}{{\c C}}1
}

\hypersetup{colorlinks=true, linkcolor=matd48azul, urlcolor=matd48azul, citecolor=matd48azul}

\newcommand{\cabecalho}[2]{%
  \begin{center}
    {\Large\bfseries\color{matd48azul} #1}\\[0.3em]
    {\large #2}\\[0.3em]
    \rule{\linewidth}{0.6pt}
  \end{center}
  \vspace{0.5em}
}

\newenvironment{questao}[1]{\vspace{1em}\noindent\textbf{Questão #1.}\hspace{0.4em}}{\par}
\newenvironment{solucao}{\vspace{0.3em}\noindent\textit{Solução.}\hspace{0.4em}\color{black!85}}{\par\vspace{0.5em}}


\begin{document}

\cabecalho{Gabarito --- Lista de Exercícios 05}{DCA com submuestreo e modelo de efeitos aleatórios (Capítulo 3 / Aula 05)}

\begin{questao}{1} \textbf{(Agricultura/pecuária --- identificação de unidades com submuestreo)}
\end{questao}
\begin{solucao}
\begin{enumerate}[label=(\alph*)]
  \item A \textbf{unidade experimental é a vaca}: é nela que a dieta é sorteada e aplicada, e não
    é possível sortear a dieta ordenha a ordenha para a mesma vaca. A \textbf{unidade amostral é
    cada ordenha}. $t=3$ dietas, $r=4$ vacas por dieta, $q=3$ ordenhas por vaca; $N = trq =
    3\times4\times3 = \textbf{36}$, que bate com $12$ vacas $\times$ 3 ordenhas.
  \item $Y_{ijk} = \mu + \tau_i + \delta_{ij} + \varepsilon_{ijk}$, $i=1,2,3$ (dieta), $j=1,\ldots,4$
    (vaca dentro da dieta $i$), $k=1,2,3$ (ordenha dentro da vaca $ij$), em que $\mu$ é a média
    geral, $\tau_i$ é o efeito fixo da dieta $i$, $\delta_{ij}\sim N(0,\sigma_\delta^2)$ é o erro
    experimental (o desvio específico da vaca $j$ na dieta $i$) e $\varepsilon_{ijk} \sim
    N(0,\sigma_\varepsilon^2)$ é o erro de amostragem (o desvio específico da ordenha $k$ daquela
    vaca), independentes entre si.
  \item Cometeria \textbf{pseudorreplicação}: trataria as $36$ ordenhas como se fossem $36$
    réplicas independentes da dieta, quando na verdade só há $12$ réplicas legítimas (vacas). As
    três ordenhas de uma mesma vaca são correlacionadas (compartilham o mesmo $\delta_{ij}$), o
    que infla artificialmente os graus de liberdade do erro usados no teste e reduz o erro-padrão
    de forma ilegítima --- tornando mais fácil rejeitar $H_0$ mesmo sem efeito real de dieta.
\end{enumerate}
\end{solucao}

\begin{questao}{2} \textbf{(Dedução --- esperança do quadrado médio do erro experimental)}
\end{questao}
\begin{solucao}
\begin{enumerate}[label=(\alph*)]
  \item Como $\bar Y_{ij\cdot} = \mu+\tau_i+\delta_{ij}+\bar\varepsilon_{ij\cdot}$ e $\bar
    Y_{i\cdot\cdot} = \mu+\tau_i+\bar\delta_{i\cdot}+\bar\varepsilon_{i\cdot\cdot}$, subtraindo os
    termos $\mu+\tau_i$ se cancelam e sobra exatamente $\bar Y_{ij\cdot}-\bar Y_{i\cdot\cdot} =
    (\delta_{ij}-\bar\delta_{i\cdot}) + (\bar\varepsilon_{ij\cdot}-\bar\varepsilon_{i\cdot\cdot})$.
    Definindo $X_{ij} := \delta_{ij}+\bar\varepsilon_{ij\cdot}$, como $\delta_{ij}$ e
    $\bar\varepsilon_{ij\cdot}$ são independentes com variâncias $\sigma_\delta^2$ e
    $\sigma_\varepsilon^2/q$ (média de $q$ termos i.i.d.), $\mathrm{Var}(X_{ij}) = \sigma_\delta^2
    + \sigma_\varepsilon^2/q$; e $\bar Y_{ij\cdot}-\bar Y_{i\cdot\cdot} = X_{ij}-\bar X_{i\cdot}$.
  \item Aplicando o fato fornecido às $r$ quantidades i.i.d. $X_{i1},\ldots,X_{ir}$ (dentro do
    tratamento $i$): $E\!\left[\sum_j (X_{ij}-\bar X_{i\cdot})^2\right] = (r-1)\,\mathrm{Var}(X_{ij})
    = (r-1)(\sigma_\delta^2+\sigma_\varepsilon^2/q)$. Somando sobre os $t$ tratamentos e
    multiplicando por $q$ (pois $SQ_{\text{exp}} = q\sum_i\sum_j(\bar Y_{ij\cdot}-\bar
    Y_{i\cdot\cdot})^2 = q\sum_i\sum_j(X_{ij}-\bar X_{i\cdot})^2$):
    $$
    E[SQ_{\text{exp}}] = q\cdot t(r-1)\left(\sigma_\delta^2+\frac{\sigma_\varepsilon^2}{q}\right)
    = t(r-1)\big(q\sigma_\delta^2+\sigma_\varepsilon^2\big).
    $$
  \item Dividindo pelos graus de liberdade $t(r-1)$ do erro experimental:
    $$
    E[MS_{\text{exp}}] = \frac{E[SQ_{\text{exp}}]}{t(r-1)} = \sigma_\varepsilon^2 +
    q\sigma_\delta^2,
    $$
    exatamente a expressão da coluna $E[QM]$ da tabela de ANOVA com submuestreo (Aula 05) para o
    erro experimental. $\blacksquare$
\end{enumerate}
\end{solucao}

\begin{questao}{3} \textbf{(Agricultura/pecuária --- ANOVA com submuestreo, tabela incompleta)}
\end{questao}
\begin{solucao}
\begin{enumerate}[label=(\alph*)]
  \item $\mathrm{g.l._{Dieta}} = t-1 = 2$; $\mathrm{g.l._{exp}} = t(r-1) = 3\times3=9$;
    $\mathrm{g.l._{am}} = tr(q-1) = 12\times2=24$ (soma $2+9+24=35$, confere).
    $SQ_{\text{Total}} = 54{,}0+39{,}0+48{,}0 = \textbf{141{,}0}$.
    $MS_{\text{Dieta}} = 54{,}0/2 = \textbf{27{,}0}$; $MS_{\text{exp}} = 39{,}0/9 \approx
    \textbf{4{,}33}$; ($MS_{\text{am}}=48{,}0/24=2{,}00$, já dado).

  \begin{center}
  \begin{tabular}{lccccc}
  \toprule
  Fonte & g.l. & SQ & QM & $F$ \\
  \midrule
  Dieta & 2 & 54{,}0 & 27{,}00 & 6{,}23 \\
  Erro experimental & 9 & 39{,}0 & 4{,}33 & --- \\
  Erro de amostragem & 24 & 48{,}0 & 2{,}00 & --- \\
  Total & 35 & 141{,}0 & & \\
  \bottomrule
  \end{tabular}
  \end{center}

  \item $F_{\text{correto}} = MS_{\text{Dieta}}/MS_{\text{exp}} = 27{,}0/4{,}33 \approx
    \textbf{6{,}23}$, com $2$ e $9$ g.l. Como $6{,}23 > F_{2,9;0{,}05}\approx 4{,}26$, rejeitamos
    $H_0$: há evidência de efeito de dieta sobre o teor de gordura.
  \item Usando o erro de amostragem por engano: $F_{\text{errado}} = MS_{\text{Dieta}}/MS_{\text{am}}
    = 27{,}0/2{,}00 = \textbf{13{,}5}$, agora com $2$ e $24$ g.l. (bem mais que o dobro do $F$
    correto, e com um denominador de graus de liberdade artificialmente maior, $24$ em vez de
    $9$). O erro é sempre nessa direção porque $MS_{\text{am}}$ estima apenas
    $\sigma_\varepsilon^2$, tipicamente menor que $MS_{\text{exp}}=\sigma_\varepsilon^2+q\sigma_\delta^2$
    --- um denominador artificialmente pequeno infla o $F$ calculado e facilita rejeitar $H_0$
    além do que os dados realmente sustentam.
\end{enumerate}
\end{solucao}

\begin{questao}{4} \textbf{(Psicologia --- componentes de variância)}
\end{questao}
\begin{solucao}
\begin{enumerate}[label=(\alph*)]
  \item \textbf{Efeitos aleatórios}: os 6 vídeos foram sorteados de uma população maior (40
    vídeos), e o interesse não está nesses 6 específicos, mas em generalizar para a população de
    vídeos possíveis --- exatamente a situação que define o modelo aleatório.
  \item $\hat\sigma_\tau^2 = (MS_{\text{Entre}}-MS_E)/n = (850-310)/10 = \textbf{54}$;
    $\hat\sigma^2 = MS_E = \textbf{310}$.
  \item $\mathrm{ICC} = \hat\sigma_\tau^2/(\hat\sigma_\tau^2+\hat\sigma^2) = 54/364 \approx
    \textbf{0{,}148}$. Cerca de 15\% da variabilidade total no tempo de resistência é atribuível à
    escolha do vídeo motivacional; os demais 85\% refletem diferenças entre participantes dentro
    de um mesmo vídeo.
\end{enumerate}
\end{solucao}

\begin{questao}{5} \textbf{(Dedução --- variância e correlação no modelo de efeitos aleatórios)}
\end{questao}
\begin{solucao}
\begin{enumerate}[label=(\alph*)]
  \item Como $\tau_i$ e $\varepsilon_{ij}$ são independentes com médias zero,
    $\mathrm{Var}(Y_{ij}) = \mathrm{Var}(\tau_i) + \mathrm{Var}(\varepsilon_{ij}) = \sigma_\tau^2 +
    \sigma^2$.
  \item Para $j\neq j'$ (mesmo nível $i$): $Y_{ij}$ e $Y_{ij'}$ compartilham o mesmo $\tau_i$, logo
    $\mathrm{Cov}(Y_{ij},Y_{ij'}) = \mathrm{Cov}(\tau_i+\varepsilon_{ij},\,\tau_i+\varepsilon_{ij'})
    = \mathrm{Var}(\tau_i) = \sigma_\tau^2$ (os termos de erro $\varepsilon_{ij}$ e
    $\varepsilon_{ij'}$ são independentes entre si e de $\tau_i$). Para $i\neq i'$ (níveis
    diferentes), $Y_{ij}$ e $Y_{i'j'}$ não compartilham nenhum termo aleatório em comum, logo
    $\mathrm{Cov}(Y_{ij},Y_{i'j'}) = 0$.
  \item $\mathrm{Corr}(Y_{ij},Y_{ij'}) = \dfrac{\mathrm{Cov}(Y_{ij},Y_{ij'})}
    {\sqrt{\mathrm{Var}(Y_{ij})\mathrm{Var}(Y_{ij'})}} = \dfrac{\sigma_\tau^2}
    {\sqrt{(\sigma_\tau^2+\sigma^2)(\sigma_\tau^2+\sigma^2)}} =
    \dfrac{\sigma_\tau^2}{\sigma_\tau^2+\sigma^2}$. Chama-se coeficiente de correlação
    \textbf{intraclasse} porque mede a correlação entre duas observações da mesma "classe"
    (mesmo nível $i$ do fator aleatório) --- e é numericamente igual à fração da variância total
    atribuível à variação \emph{entre} classes, tal como calculado na Questão 4(c).
\end{enumerate}
\end{solucao}

\end{document}
